\section{Probability Model: The Foundation of Inference}

\subsection{Descriptive vs. Inferential Statistics}

In Data Science, we deal with two fundamentally different problems. It is crucial to distinguish them immediately.

\begin{defbox}{Descriptive Statistics (The Direct Problem)}
We observe the \textbf{entire population} (e.g., every single customer in a database).
\begin{itemize}
    \item \textbf{Goal:} Summarize the data exactly.
    \item \textbf{Result:} Certain facts (e.g., \textit{The exact average age is 35.2}).
    \item \textbf{Uncertainty:} Zero.
\end{itemize}
\end{defbox}

\begin{defbox}{Statistical Inference (The Indirect Problem)}
We only observe a small \textbf{sample} (e.g., 100 surveyed customers).
\begin{itemize}
    \item \textbf{Goal:} Infer the characteristics of the whole population based on the sample.
    \item \textbf{Result:} Estimates with uncertainty (e.g., \textit{The avg age is likely between 34 - 36}).
    \item \textbf{Tool:} We use Probability to measure that uncertainty.
\end{itemize}
\end{defbox}

\paragraph{Basic Definitions. }

To move from a sample to a population, we must formally define the entities involved and the goals of inference.
\begin{itemize}[label=$\cdot$]
    \item \textbf{Population}: the entire group under study to which conclusions are to be applied.
    \item \textbf{Sample}: a subset of the population that is actually observed and collected.
    \item \textbf{Statistical Inference}: The process of drawing conclusions about the population from data collected on a sample, using probability reasoning.
\end{itemize}

\subsubsection{The Two Main Tasks of Inference}
As a Data Scientist, you will mostly perform two tasks:
\begin{enumerate}
\item \textbf{Estimation:}  This involves estimating unknown population characteristics (parameters) based on sample statistics.
 \begin{itemize}
  \item \textbf{Point Estimation}: Assigning a single value to the population parameter (e.g., using the sample mean $\bar{x}$ to estimate the population mean $\mu$).
  \item \textbf{Interval Estimation}: Assigning an interval of plausible values (a confidence interval) to the parameter.
  \end{itemize}
    \item \textbf{Hypothesis Testing:} Deciding whether a claim (hypothesis) about a population parameter is acceptable or should be rejected, in probabilistic terms.(e.g., Did the new marketing campaign actually increase sales or was it just luck?).
\end{enumerate}

\vspace{0.2cm}
\begin{example}{\textbf{Parameter Estimation}}.

A Data Scientist wants to know the average age ($\mu$) of all customers of an e-commerce platform (Population). Since they cannot survey millions of users, they take a sample of $n=100$ customers.
\begin{itemize}
 \item \texttt{Step1 (Descriptive)}:
 \begin{itemize}
     \item Calculate the sample mean age. Suppose the sample mean is $\bar{x} = 35.2$ years.
 \end{itemize}
 \item \texttt{Step2 (Point Estimation)}:
 \begin{itemize}
     \item The point estimate for the population mean $\mu$ is $35.2$.
 \end{itemize}
 \item \texttt{Step3 (Interval Estimation)}:
 \begin{itemize}
     \item Based on sample var and size, they calculate a 95\% Confidence Interval for $\mu$, say $[34.0, 36.4]$. They infer that the true $\mu$ age of \textit{all} customers is likely between 34.0 and 36.4 y.o.
 \end{itemize} 
\end{itemize}
\end{example}

\subsubsection{The Importance of Unbiased Random Samples}

Since observations are made only on a sample, the fundamental questions are:
\begin{itemize}[label=-]
 \item How to extend sample conclusions to the population?
 \item What are the limits of validity and reliability of these conclusions?
\end{itemize}

The solution lies in using an \textbf{Unbiased Random Sample}:

\begin{itemize}
\item  A sample where the units are selected through a \textbf{random draw}, ensuring that the selection process is not influenced by the studied variable. This means every unit in the population has the same probability of being chosen.
\begin{itemize}
    \item This method is the best approach to \textbf{ensure representativeness and validity of inference}.
    \item If the survey is not random, inferential methods cannot be applied because the sample might not accurately reflect the population.
\end{itemize}
\end{itemize}

\subsection{Probability: The Measure of Uncertainty}

\subsubsection{Why Do We Need Probability?}

In the previous section, we learned that we often only have a \textbf{Sample}, but we want to know about the \textbf{Population}.
Because the sample is just a small piece of the whole, our conclusions are never 100\% certain.

\textbf{Probability} is the mathematical tool we use to measure that uncertainty. It allows us to say: \textit{I am 95\% confident that the average age is between 34 and 36,} rather than just guessing.

\subsubsection{The Three Definitions of Probability}

How do we actually define \textit{Chance}? There are three main schools of thought. As a Data Scientist, you will mostly use the \textbf{Classical} and \textbf{Frequentist} approaches, but \textbf{Subjectivist} is the basis of Bayesian Statistics.

\begin{enumerate}
    \item \textbf{Classical Probability (Gambling):}
    Used when all outcomes are equally likely (e.g., dice, cards).
    \[ P(E) = \frac{\text{Number of Favorable Outcomes}}{\text{Total Number of Possible Outcomes}} \]

\vs
    \item \textbf{Frequentist Probability (Simulation):}
    Used for repeatable experiments. Probability is the \textbf{relative frequency} of an event happening over a huge number of trials.
    \[ P(E) \approx \frac{\text{Count of Successes}}{\text{Count of Trials}} \quad (\text{as trials } \to \infty) \]

\vs
    \item \textbf{Subjectivist Probability (Belief):}
    Used when experiments cannot be repeated (e.g., \textit{Will humanity go to Mars by 2030?}). It measures a degree of personal belief or betting odds.
\end{enumerate}

\begin{example}{\textbf{Example: Tossing Coins}}

If you toss two coins, the Sample Space is $\{HH, HT, TH, TT\}$.
The event \texttt{Getting exactly one head} has 2 favorable outcomes ($HT, TH$).
\[ P(\text{One Head}) = \frac{2}{4} = 0.5 \]
\end{example}

\pagebreak
\subsection{Axioms and Fundamental Rules}

\subsubsection{Set Theory and Venn Diagrams}
We treat Probability mathematically using \textbf{Set Theory}. We visualize outcomes using \textbf{Venn Diagrams}.

\begin{itemize}[label=$\cdot$]
\item \textbf{Trial}: is a single execution of an experiment or process that results in an outcome (e.g.,The single roll of the die).

\vs
    \item \textbf{Sample Space ($S$):} the set of all possible outcomes of an experiment. The rectangle containing everything that can happen.

    \vs
    
    \item \textbf{Event ($A$):} a collection of outcomes, or a subset of the sample space S. A circle inside the rectangle representing a specific outcome. Events are combined using Boolean Logic Operators:  Union $\cup$, Intersection $\cap$,  Complement $\overline{A}$.
\end{itemize}


% IMMAGINE ORIGINALE MANTENUTA
\begin{figure}[H] 
    \centering\includegraphics[scale=0.4]{103.Venn Diagram.jpeg}
    \caption{Venn Diagram}
    \captionsetup{font=footnotesize}
    \label{fig-Venn Diagram}
\end{figure}

\subsubsection{Kolmogorov's Three Axioms}
The \textbf{three fundamental axioms} define the mathematical properties that any probability measure $P$ must satisfy.

\begin{enumerate}
    \item \textbf{Positivity:} You cannot have a negative chance.
    \[ P(E) \ge 0 \]
    \item \textbf{Certainty:} The probability that \textit{something} in the sample space happens is 100\%.
    \[ P(S) = 1 \]
    \item \textbf{Additivity (for Disjoint Events):} If two events cannot happen at the same time (Mutually Exclusive, $A \cap B = \emptyset$), you just add their probabilities.
    \[ P(A \cup B) = P(A) + P(B) \]
\end{enumerate}

\subsubsection{Important Theorems derived from Axioms}

From those three simple rules, we derive the formulas you will use in practice.

\begin{itemize}
 \item \textbf{The Complement Rule}: The probability that an event $E$ does \textit{not} occur ($\overline{E}$) is one minus the probability that $E$ occurs. This implies that $0 \le P(E) \le 1$.
  \begin{equation}
      P(\overline{E}) = 1 - P(E)
  \end{equation}

    \item \textbf{The General Addition Law (Union):}
    If events $A$ and $B$ \textbf{can} overlap (happen together), we cannot just add them, or we would double-count the overlap. We must subtract the intersection.
    \[ P(A \text{ or } B) = P(A) + P(B) - P(A \text{ and } B) \]



 \item \textbf{Probability of the Empty Set ($\emptyset$)}: The probability of the impossible event is zero. This is derived from Axiom 3 (Additivity) and Axiom 2 (Unit Measure).
  \begin{equation}
      P(\emptyset) = 0
  \end{equation}
\end{itemize}


% IMMAGINE ORIGINALE MANTENUTA
\begin{figure}[H] 
    \centering\includegraphics[scale=0.45]{104.Fundamental Theorems.jpeg}
    \caption{General Addition Law}
    \captionsetup{font=footnotesize}
    \label{fig-Fundamental Theorems}
\end{figure}

\begin{example}{\textbf{Example: Passing Exams}}

Let $P(Stat) = 0.60$ and $P(CS) = 0.50$.
If 35\% of students pass \textbf{both} ($P(Stat \cap CS) = 0.35$), what is the chance of passing \textbf{at least one}?
\[ P(Stat \cup CS) = 0.60 + 0.50 - 0.35 = 0.75 \]
\end{example}

\subsubsection{Extension: Three Events}
If you have three overlapping events ($A, B, C$), the logic is the same: Add the single circles, subtract the pairwise overlaps, and add back the center (because you subtracted it too many times).
\[ P(A \cup B \cup C) = P(A) + P(B) + P(C) - P(A \cap B) - P(A \cap C) - P(B \cap C) + P(A \cap B \cap C) \]

\subsection{Conditional Probability}

This is the most critical concept for understanding how Machine Learning models (like Naive Bayes) work. It answers: \textit{How do my chances change if I get new information?}

\subsubsection{Defining Conditional Probability}

\begin{defbox}{Conditional Probability}
The probability of event $B$ happening, \textbf{given that} we already know event $A$ has happened.
\begin{equation}
    P(B|A) = \frac{P(A \cap B)}{P(A)}
\end{equation}
\end{defbox}
We shrink the sample space. We no longer care about the whole world ($S$), only the world where $A$ is true.


\begin{defbox}{The Factorization Theorem (Total Probability)}
If we want to find the total probability of $B$, we sum up all the ways $B$ can happen (via $A$ or via $\bar{A}$):
\begin{equation}
    P(B) = P(A \cap B) + P(\bar{A} \cap B) =P(B|A)P(A) + P(B|\bar{A})P(\bar{A})
\end{equation}
\end{defbox}

\subsubsection{Stochastic Independence}
Two events are \textbf{Independent} if knowing $A$ gives us zero information about $B$.
\begin{equation}
    P(B|A) = P(B)
\end{equation}
If (and only if) this holds, we can just multiply probabilities: $P(A \cap B) = P(A) \cdot P(B)$.

\vspace{0.2cm}

\begin{example}{\textbf{Example 1: The Two-Way Table}}

Let's apply this to a real dataset.
$A$ = Female, $B$ = Wears Glasses.

% IMMAGINE ORIGINALE MANTENUTA
\begin{figure}[H] 
    \centering\includegraphics[scale=0.45]{136.Probability Two way table.jpeg}
    \caption{Probability Two way table}
    \captionsetup{font=footnotesize}
    \label{fig-Probability Two way table}
\end{figure}

\begin{itemize}
    \item \textbf{Question:} What is the probability a person wears glasses ($B$), \textit{given} they are Female ($A$)?
    \item \textbf{Calculation:} Look only at the \textit{Female} row.
    \[ P(B|A) = \frac{20 \text{ (Females with glasses)}}{40 \text{ (Total Females)}} = 0.5 \]
    \item \textbf{Check Independence:}
    $P(B) = 60/100 = 0.6$. Since $P(B|A) \neq P(B)$ ($0.5 \neq 0.6$), Gender and Glasses are \textbf{Dependent}.
\end{itemize}
\end{example}

\subsubsection{Bayes' Theorem}

This is the most important formula in Data Science. It tells us how to \textit{update our beliefs} when we see new evidence.

\begin{defbox}{Bayes' Theorem}
\begin{equation}
    P(A|E) = \frac{P(E|A) \cdot P(A)}{P(E)}
\end{equation}
\begin{itemize}
    \item $P(A)$: \textbf{Prior}. What we believed \textit{before} seeing data.
    \item $P(E|A)$: \textbf{Likelihood}. How probable is the evidence if our hypothesis is true?
    \item $P(A|E)$: \textbf{Posterior}. Our new belief \textit{after} seeing the evidence.
\end{itemize}
\end{defbox}

\subsubsection{Real World Application: Medical Diagnosis}
This is a classic interview question to test if you understand \textit{Base Rate Fallacy}.

\begin{example}{\textbf{The False Positive Paradox}}
\textbf{Scenario:}
\begin{itemize}
    \item A rare disease affects 1\% of people ($P(D) = 0.01$).
    \item The test is 95\% accurate ($P(+|D) = 0.95$).
    \item The test has a 10\% false positive rate ($P(+|\text{Health}) = 0.10$).
\end{itemize}
\textbf{Question:} You test positive. What is the probability you actually have the disease?

\textbf{The Calculation (Bayes):}
\[ P(D|+) = \frac{P(+|D) \cdot P(D)}{P(+)} \]

\begin{enumerate}
    \item \textbf{Numerator (True Positive path):} $0.95 \times 0.01 = 0.0095$
    \item \textbf{Denominator (Total Positives):}
    True Positives + False Positives
    \[ (0.95 \times 0.01) + (0.10 \times 0.99) = 0.0095 + 0.099 = 0.1085 \]
    \item \textbf{Result:}
    \[ \frac{0.0095}{0.1085} \approx \textbf{8.76\%} \]
\end{enumerate}

\textbf{Conclusion:} Even with a positive test, you are likely healthy! Why? Because the disease is so rare ($1\%$) that the false alarms ($10\%$ of the $99\%$ healthy people) drown out the real cases.
\end{example}



\subsection{Combinatorial Calculus and Discrete Random Variables}

\subsubsection{Combinatorial Calculus}

Before calculating probabilities, we often need to count how many ways something can happen. This is \textbf{Combinatorics}.

\begin{defbox}{Ordered Sequences (Does Order Matter? YES)}
\textbf{1. With Replacement:}
You have $n$ options, and you choose $k$ times. After each choice, you put the item back.
\begin{equation}
n^k
\end{equation}
\begin{itemize}
    \item \textit{Example:} A 4-digit PIN code ($10^4 = 10,000$ codes).
\end{itemize}

\vs
\textbf{2. Without Replacement (Permutations):}
You have $n$ options, choose $k$. Once picked, an item is gone
\begin{equation}
D_{n,k} = \frac{n!}{(n-k)!}
\end{equation}
\begin{itemize}
    \item \textit{Example:} A race with 8 runners. Who gets Gold, Silver, and Bronze ($k=3$)?

\end{itemize}
\end{defbox}

\begin{defbox}{Unordered Sequences (Does Order Matter? NO)}
\textbf{Combinations (The Binomial Coefficient):}
You choose $k$ items from $n$, and the order does \textbf{not} matter.
    \begin{equation}
        C_{n,k} = \binom{n}{k} = \frac{n!}{k!(n-k)!}
    \end{equation}
\begin{itemize}
    \item \textit{Example:} Choosing a team of 3 Data Scientists from a department of 8.
\end{itemize}
\end{defbox}

This is the most critical formula for Data Science (used in sampling).

\subsubsection{Discrete Random Variables}

A \textbf{Random Variable} ($X$) is just a function that translates a \textit{real world outcome} into a number. It is \textbf{Discrete} because can only take specific integer values (e.g., number of sales, number of clicks).


\pagebreak
\begin{example}{\textbf{Discrete Random Variable}}

  Imagine an experiment, like flipping a coin three times. The raw outcomes (elementary events) are: {HHH, HHT, HTH, THH, HTT, THT, TTH, TTT}.
  
  \vs
  The Discrete Random Variable ($X$) is defined as the "Number of Heads" in the three flips: 3, 2, 2, ...  
\end{example}

         \begin{figure}[H] 
        \centering\includegraphics[scale=0.15]{106.Discrete variable.jpeg}
        \caption{Discrete Random variable}
        \captionsetup{font=footnotesize}
        \label{fig-Discrete Random variable}
    \end{figure}


\begin{defbox}{Key Metrics: Mean and Variance}
\textbf{1. Expected Value (mean):}
If you repeated the experiment infinitely, what would be the average?
\begin{equation}
    E(X) = \mu = \sum x_i \cdot p_i
\end{equation}

\textbf{2. Variance:}
How far are the values typically from the mean?
\begin{equation}
    \text{Var}(X) = E[(X - \mu)^2]
\end{equation}
\end{defbox}


\subsubsection{Discrete Probability Models}

In the previous section, we learned what a Random Variable is (a way to map outcomes to numbers). Now we'll explore several important \textbf{probability distributions} that model common real-world scenarios.

Think of these as  \textit{patterns} that describe how data behaves in the real world. Once you recognize the pattern (e.g., \textit{counting successes}), you can just pick the right tool (e.g., \textit{Binomial}) and use its ready-made formulas.

\vspace{0.3cm}
\textbf{1. The Basic Building Blocks}


\vspace{0.05cm}

\begin{defbox}{Uniform Distribution $U(k)$}
\textbf{Scenario:} \textit{Everything is equally likely.}
\begin{itemize}
    \item \textbf{Example:} Rolling a fair die. k is the number of possible outcomes.
    \item \textbf{Probability:} 
    \begin{equation}P(X=x) = \frac{1}{k}\end{equation}
   
    \item \textbf{Mean:} The exact middle value. 
\end{itemize}
\end{defbox}

\begin{figure}[H] 
    \centering\includegraphics[scale=0.13]{201.Discrete uniform distribution.jpeg}
    \caption{Discrete uniform distribution}
    \captionsetup{font=footnotesize}
\end{figure}




\begin{defbox}{Bernoulli Distribution $B(1, \pi)$}
\textbf{Scenario:} \textit{The Atom of Probability.} A single trial with only two outcomes: Success (1) or Failure (0).
\begin{itemize}
    \item \textbf{Example:} A user clicks an ad (Yes/No).
    \item \textbf{Parameter $\pi$:} The probability of success.
    \item \textbf{Mean:} $\pi$.
    \item \textbf{Variance:} $\pi(1-\pi)$ (Variance is highest when $\pi=0.5$: maximum uncertainty).
\end{itemize}
\end{defbox}

\begin{figure}[H] 
    \centering\includegraphics[scale=0.18]{202.Bernoulli Distribution.jpeg}
    \caption{Bernoulli Distribution}
    \captionsetup{font=footnotesize}
\end{figure}


\paragraph*{Why is the mean of a Bernoulli variable equal to \(\pi\)?}
A Bernoulli random variable \(X\) can take only two values:
\[
X =
\begin{cases}
1 & \text{if success occurs} \\
0 & \text{if failure occurs}
\end{cases}
\]

The parameter \(\pi\) is the probability of success, that is:
\[
P(X=1)=\pi, \qquad P(X=0)=1-\pi
\]

The mean of a random variable is its expected value, computed as a weighted average of its possible numerical values:
\[
E[X] = 1 \cdot P(X=1) + 0 \cdot P(X=0)
\]

Substituting the Bernoulli probabilities, we get:
\[
E[X] = 1 \cdot \pi + 0 \cdot (1-\pi) = \pi
\]

So, the mean is equal to \(\pi\) because the variable is coded as \(1\) for success and \(0\) for failure. In other words, the average value of a Bernoulli variable represents the proportion of successes over many repeated trials.

\paragraph*{Why is the variance of a Bernoulli variable equal to \(\pi(1-\pi)\)?}
A Bernoulli random variable \(X\) takes only two possible values:
\[
X =
\begin{cases}
1 & \text{with probability } \pi \\
0 & \text{with probability } 1-\pi
\end{cases}
\]

We already know that its mean is:
\[
E[X]=\pi
\]

To compute the variance, we use:
\[
\mathrm{Var}(X)=E[X^2]-(E[X])^2
\]

For a Bernoulli variable, \(X^2=X\), because:
\[
0^2=0, \qquad 1^2=1
\]

Therefore,
\[
E[X^2]=E[X]=\pi
\]

Substituting into the variance formula:
\[
\mathrm{Var}(X)=\pi-\pi^2=\pi(1-\pi)
\]

This result shows that the variance depends on how balanced success and failure are. It is small when \(\pi\) is close to 0 or 1, because the outcome is more predictable. It is maximum when \(\pi=0.5\), because success and failure are equally likely, giving the highest uncertainty.


\vspace{0.25cm}
\textbf{2. Counting Successes (The \textit{Binomial} Family)}

These distributions answer the question: \textit{How many successes did I get?}

\vs

\begin{defbox}{Binomial Distribution $B(n, \pi)$}
\textbf{Scenario:} You repeat the same Bernoulli trial \(n\) times independently, and you count how many successes occur.

\begin{itemize}
    \item \textbf{Example:} Suppose you perform \(12\) independent trials (\(n=12\)), and each trial has probability \(0.4\) of success (\(\pi=0.4\)). Let \(X\) be the number of successes. The graph shows the probability of obtaining exactly \(0,1,2,\dots,12\) successes.

\vs

    \item \textbf{Probability:}
    \begin{equation}
    P(X=x)=\binom{n}{x}\pi^x(1-\pi)^{n-x}, \qquad x=0,1,\dots,n
    \end{equation}
    \item \textbf{Mean:}
    \begin{equation}
    E[X]=n\pi=12 \cdot 0.4=4.8
   \end{equation}
    This explains why the distribution is centered around 5 successes.
\end{itemize}
\end{defbox}

\begin{figure}[H] 
    \centering\includegraphics[scale=0.23]{203.Binomial Distribution.jpeg}
    \caption{Binomial Distribution}
    \captionsetup{font=footnotesize}
\end{figure}

Each bar represents \(P(X=x)\), that is, the probability of getting exactly \(x\) successes. In this example, the tallest bars are around \(x=4\) and \(x=5\), meaning that getting about 4 or 5 successes is the most likely outcome.

\vs

\begin{defbox}{Hypergeometric Distribution $H(N, n, m)$}
\textbf{Scenario:} Similar to the Binomial case, but now the sampling is done \textbf{without replacement}, so the trials are dependent.

\begin{itemize}
    \item \textbf{Example:} Suppose a population contains \(N=30\) items, of which \(m=10\) are classified as successes. You randomly draw \(n=8\) items without replacement. Let \(X\) be the number of successes in the sample.
    \vs

    \item \textbf{Key Difference from Binomial:} The trials are \textbf{dependent}, because the composition of the population changes after each draw. If one success is drawn early, there is one fewer success left in the population, so the probability changes at the next draw.
    \vs

\item \textbf{Probability:}
    \begin{equation}
P(X=x)=\frac{\binom{m}{x}\binom{N-m}{\,n-x\,}}{\binom{N}{n}}
   \end{equation}
where:
\begin{itemize}
    \item \(N\) = total number of items in the population.
    \item \(m\) = number of success states in the population.
    \item \(n\) = number of draws.
    \item \(x\) = number of observed successes in the sample.
\end{itemize}
    \vs

\item \textbf{Mean:}
\[
E[X]=n\frac{m}{N}
\]
This means that the expected number of successes in the sample is equal to the sample size multiplied by the proportion of successes in the population.

\end{itemize}
\end{defbox}

\begin{figure}[H] 
    \centering\includegraphics[scale=0.23]{204.Hypergeometric Distribution.jpeg}
    \caption{Hypergeometric Distribution}
    \captionsetup{font=footnotesize}
\end{figure}

Each bar represents \(P(X=x)\), the probability of obtaining exactly \(x\) successes in the sample. In this example, the highest probabilities are around \(x=2\) and \(x=3\), meaning that getting 2 or 3 successes is the most likely outcome.

\textbf{Usage:} This model is useful when sampling from a finite population, especially in applications such as quality control, auditing, and card-drawing problems.

\vspace{0.3cm}

\textbf{3. Waiting Time (The \textit{Geometric} Family)}

These distributions answer the question: \textit{How long do I have to wait?}

\begin{defbox}{Geometric Distribution \(G(\pi)\)}
\textbf{Scenario:} The Geometric distribution models the number of trials needed until the \textbf{first success} occurs.

\begin{itemize}
    \item \textbf{Example:} Suppose each trial has probability \(p=0.25\) of success, and let \(X\) be the trial on which the first success occurs. The graph shows the probability of getting the first success on trial \(1,2,3,\dots\).

   
    \item \textbf{Formula:}
    \begin{equation}
    P(X=x)=(1-p)^{x-1}p, \qquad x=1,2,3,\dots
    \end{equation}

    \item \textbf{Mean:}
    \begin{equation}
    E[X]=\frac{1}{p}
  \end{equation}
    For example, if \(p=0.25\), then
    \[
    E[X]=\frac{1}{0.25}=4
    \]
    so on average we expect the first success around the 4th trial.

    \item \textbf{Property:} \textit{Memoryless}. If no success has occurred yet, the probability of success on the next trial is still \(p\), exactly the same as at the beginning.
\end{itemize}
\end{defbox}

\begin{figure}[H] 
    \centering\includegraphics[scale=0.23]{205.Geometric Distribution.jpeg}
    \caption{Geometric Distribution}
    \captionsetup{font=footnotesize}
\end{figure}

Each bar represents \(P(X=x)\), the probability that the first success occurs exactly on trial \(x\). The highest probability is at \(x=1\), and then the probabilities decrease as \(x\) increases. This means that early success is more likely, while having to wait many trials is possible but less likely.


\vs
\begin{defbox}{Negative Binomial \(NB(k,\pi)\)}
\textbf{Scenario:} The Negative Binomial distribution models the number of trials needed until the \textbf{\(k\)-th success} occurs.

\begin{itemize}
    \item \textbf{Example:} Suppose each trial has probability \(p=0.35\) of success, and let \(X\) be the trial on which the \(4\)-th success occurs. The graph shows the probability that the 4-th success happens on trial \(4,5,6,\dots\).

   

    \item \textbf{Formula:}
    \begin{equation}
    P(X=x)=\binom{x-1}{k-1}p^k(1-p)^{x-k}, \qquad x=k,k+1,k+2,\dots
    \end{equation}

    \item \textbf{Mean:}
     \begin{equation}
    E[X]=\frac{k}{p}
    \end{equation}
    For example, if \(k=4\) and \(p=0.35\), then
    \[
    E[X]=\frac{4}{0.35}\approx 11.43
    \]
    so on average we expect the 4-th success around the 11-th trial.

\end{itemize}
\end{defbox}

\begin{figure}[H] 
    \centering\includegraphics[scale=0.23]{206.negative binomial.jpeg}
    \caption{Negative Binomial Distribution}
    \captionsetup{font=footnotesize}
\end{figure}



Each bar represents \(P(X=x)\), the probability that the \(k\)-th success occurs exactly on trial \(x\). In this example, the highest probabilities are around \(x=9\) and \(x=10\), meaning that the 4-th success is most likely to occur around the 9-th or 10-th trial.

\textbf{Connection with the Geometric distribution:} It can be seen as the sum of \(k\) independent Geometric waiting times, because we are waiting for the first success, then the second, and so on, until the \(k\)-th success.

\vspace{0.3cm}
\textbf{4. Rate of Events (Poisson)}

\vs
\begin{defbox}{Poisson Distribution \(P(\lambda)\)}
\textbf{Scenario:} The Poisson distribution models the number of events occurring in a fixed interval of time or space, assuming that events happen independently and at an average rate \(\lambda\).

\begin{itemize}
    \item \textbf{Example:} Suppose the average number of events is \(\lambda=4\) in a given interval, and let \(X\) be the number of events observed in that interval. The graph shows the probability of observing exactly \(0,1,2,\dots\) events.


    \item \textbf{Formula:}
    \begin{equation}
    P(X=x)=\frac{e^{-\lambda}\lambda^x}{x!}, \qquad x=0,1,2,\dots
    \end{equation}

    \item \textbf{Parameter \(\lambda\):} \(\lambda\) is the average number of events in the interval. It determines both the center and the spread of the distribution.

    \item \textbf{Mean and Variance:}
    \begin{equation}
    E[X]=\lambda, \qquad \mathrm{Var}(X)=\lambda
  \end{equation}
    For example, if \(\lambda=4\), then the expected number of events is 4, and the variance is also 4.


\end{itemize}
\end{defbox}

\begin{figure}[H] 
    \centering\includegraphics[scale=0.25]{207.Poisson Distribution.jpeg}
    \caption{Poisson Distribution}
    \captionsetup{font=footnotesize}
\end{figure}

 Each bar represents \(P(X=x)\), the probability of observing exactly \(x\) events. In this example, the highest probabilities are around \(x=3\) and \(x=4\), meaning that observing 3 or 4 events is the most likely outcome.

    \textbf{Usage:} This model is widely used for counting events such as arrivals, failures, calls, or accidents over a fixed interval.

% IMMAGINE ORIGINALE MANTENUTA
\begin{figure}[H] 
    \centering\includegraphics[scale=0.5]{107.Discrete probabilty model tabel.jpeg}
    \caption{Discrete probabilty model tabel}
    \captionsetup{font=footnotesize}
    \label{fig-Discrete probabilty model tabel}
\end{figure}

\pagebreak

\subsection{Continuous Random Variables and Probability Models}

\subsubsection{Continuous Random Variables (PDF \& CDF)}
In the continuous world, variables can take infinite possible decimal values. The probability of picking one \textit{exact} number is mathematically zero. Instead, we measure the probability of a value falling within a specific interval by calculating the \textbf{area under a curve}.

\begin{itemize}
    \item \textbf{Probability Density Function (PDF)  $f(x)$:} Represents the \textit{shape} of the distribution. It must always be positive, and the total area under the curve must equal 1 ($100\%$ probability). It tells you the height of the curve at the point $x$. 
    \item \textbf{Cumulative Distribution Function (CDF)  $F(x)$:} Represents the \textit{running total} of probability. It is the area under the curve from $-\infty$ up to a specific point $x$.
    \begin{equation}
        F(x) = P(X \leq x) = \int_{-\infty}^{x} f(t) \, dt
    \end{equation}
\end{itemize}

\textbf{The Moments} \\
The logic for mean and variance is identical to discrete variables, but we replace the summation symbol ($\sum$) with an integral ($\int$):
\begin{itemize}
    \item \textbf{Expected Value (Mean):} $E(X) = \int_{-\infty}^{+\infty} x f(x) \, dx = \mu$
    \item \textbf{Variance:} $\text{Var}(X) = \int_{-\infty}^{+\infty} (x-\mu)^2 f(x) \, dx = \sigma^2$
\end{itemize}

\begin{figure}[H] 
    \centering\includegraphics[scale=0.15]{108.Continuous Random Variable.jpeg}
    \caption{Continuous Random Variable}
    \captionsetup{font=footnotesize}
    \label{fig-Continuous Random Variable}
\end{figure}

This figure does not represent the CDF directly. It shows the probability over a finite interval, that is, the area under the density curve between \(x_0\) and \(x_1\). By contrast, the CDF at a point \(x\) is the total area under the curve from \(-\infty\) up to \(x\).

\subsubsection{Key Continuous Probability Models}
As a Data Scientist, you will encounter these foundational distributions repeatedly depending on the type of data you are modeling:

\begin{defbox}{Uniform Distribution \(X \sim U_c(a,b)\)}
\textbf{Scenario:} A continuous random variable is uniformly distributed over the interval \([a,b]\), meaning that all values inside the interval are equally likely in terms of density.

\begin{itemize}
 \item \textbf{Probability density function:}The density function is constant over \([a,b]\), so the graph is a horizontal line. Its height is
   \item \textbf{Probability density function:}
    \begin{equation}
    f(x)=
    \begin{cases}
    \dfrac{1}{b-a}, & a \le x \le b \\[6pt]
    0, & \text{otherwise}
    \end{cases}
\end{equation}

    for \(a \le x \le b\), and \(f(x)=0\) outside the interval.


\end{itemize}

\end{defbox}

   Since the density is flat, intervals of the same length inside \([a,b]\) have the same probability.


\begin{figure}[H] 
    \centering\includegraphics[scale=0.25]{137.Uniform distribution.jpeg}
    \caption{Uniform distribution}
    \captionsetup{font=footnotesize}
\end{figure}

 The density forms a rectangle with base \(b-a\) and height \(\frac{1}{b-a}\), so the total area is
    \[
    (b-a)\cdot \frac{1}{b-a}=1
    \]
    as required for a probability density function.

\begin{defbox}{ Exponential Distributions }



  The Exponential distribution models the waiting time until the \textbf{next event} occurs. It is commonly used when events happen independently at a constant average rate.


\begin{itemize}
  

    \item \textbf{Probability density function:}
    \begin{equation}
    f(x)=
    \begin{cases}
    \lambda e^{-\lambda x}, & x \geq 0 \\[4pt]
    0, & x<0
    \end{cases}
\end{equation}

    \item \textbf{Mean:}
    \begin{equation}
    E[X]=\frac{1}{\lambda}
 \end{equation}

    \item \textbf{Property:} The Exponential distribution is \textit{memoryless}. This means that if no event has occurred yet, the waiting time for the next event does not depend on how long we have already waited.
    
\end{itemize}
\end{defbox}

\begin{figure}[H] 
    \centering\includegraphics[scale=0.2]{138.Esponential Distribution.jpeg}
    \caption{Esponential Distribution}
    \captionsetup{font=footnotesize}
\end{figure}



The density is highest near \(x=0\) and then decreases continuously as \(x\) increases. This means that short waiting times are more likely, while very long waiting times are possible but less likely.

\begin{defbox}{Gamma Distributions $(\kappa,\lambda)$}

The Gamma distribution is a generalization of the Exponential distribution. It models the waiting time until \(\kappa\) events occur.


    \begin{itemize}
 

    \item \textbf{Probability density function:}
    \begin{equation}
    f(x)=
    \begin{cases}
    \dfrac{\lambda^\kappa x^{\kappa-1}e^{-\lambda x}}{\Gamma(\kappa)}, & x \geq 0 \\[8pt]
    0, & x<0
    \end{cases}
\end{equation}

    \item \textbf{Mean:}
  \begin{equation}
    E[X]=\frac{\kappa}{\lambda}
  \end{equation}

\end{itemize}

\end{defbox}

\begin{figure}[H] 
    \centering\includegraphics[scale=0.2]{139.Gamma Distribution.jpeg}
    \caption{Gamma Distribution}
    \captionsetup{font=footnotesize}
\end{figure}


Unlike the Exponential case, the Gamma density may first increase, reach a peak, and then decrease. This reflects the fact that we are no longer waiting for just one event, but for the \(\kappa\)-th event.

\textbf{Connection with the Exponential distribution:} The Exponential distribution is a special case of the Gamma distribution when \(\kappa=1\).


\begin{defbox}{Beta Distribution \(X \sim \mathrm{Be}(\alpha,\beta)\)}
The Beta distribution models a continuous random variable constrained to the interval \([0,1]\). It is very flexible and can take many different shapes depending on the values of \(\alpha\) and \(\beta\).

\begin{itemize}


    \item \textbf{Probability density function:}
    \begin{equation}
    f(x)=
    \begin{cases}
    \dfrac{1}{B(\alpha,\beta)}x^{\alpha-1}(1-x)^{\beta-1}, & 0 \le x \le 1 \\[8pt]
    0, & \text{otherwise}
    \end{cases}
\end{equation}
    where \(B(\alpha,\beta)\) is the Beta function.

    \item \textbf{Mean:}
    \begin{equation}
    E[X]=\frac{\alpha}{\alpha+\beta}
  \end{equation}

  \vs

    \item \textbf{Why it is useful:} Since the Beta distribution is bounded between 0 and 1, it is especially useful for modeling probabilities, proportions, and rates.


\end{itemize}
\end{defbox}

\begin{figure}[H] 
    \centering\includegraphics[scale=0.2]{140.beta distribution.jpeg}
    \caption{Beta Distribution}
    \captionsetup{font=footnotesize}
\end{figure}

The density is defined only between 0 and 1. In this example, the curve is concentrated more toward smaller values of \(x\), meaning that values closer to 0 are more likely than values close to 1.

 \textbf{Data Science Use Case:} It is widely used in Bayesian statistics to model uncertain probabilities, such as conversion rates, click-through rates, and success proportions.



\subsubsection{The Normal (Gaussian) Distribution}
\begin{defbox}{Normal Distribution \(N(\mu,\sigma^2)\)}
The Normal distribution is one of the most important distributions in statistics, mainly because of the \textbf{Central Limit Theorem}. It is a continuous distribution with a characteristic bell-shaped curve, perfectly symmetric around its mean \(\mu\).

\begin{itemize}
    \item \textbf{Parameters:}
    \[
    \mu = \text{mean}, \qquad \sigma^2 = \text{variance}
    \]
    The parameter \(\mu\) determines the center of the distribution, while \(\sigma^2\) controls the spread. A larger \(\sigma\) produces a wider and flatter curve.

    \vs
    \item \textbf{Probability density function:}
    \begin{equation}
    f(x)=\frac{1}{\sqrt{2\pi\sigma^2}} \exp\left(-\frac{(x-\mu)^2}{2\sigma^2}\right), \qquad x\in\mathbb{R}
    \end{equation}




\end{itemize}
\end{defbox}

\begin{figure}[H] 
    \centering\includegraphics[scale=0.45]{110.Normal Distribution.jpeg}
    \caption{Normal Distribution showing the empirical rule}
    \captionsetup{font=footnotesize}
    \label{fig-Normal Distribution}
\end{figure}

\textbf{Symmetry:} Since the curve is perfectly symmetric around \(\mu\), the mean, median, and mode all coincide:
    \[
    \text{mean}=\text{median}=\text{mode}=\mu
    \]

  \textbf{Empirical Rule (68--95--99.7):}
    \begin{itemize}
        \item Approximately \(68\%\) of the data falls within \(1\sigma\) of the mean.
        \item Approximately \(95\%\) of the data falls within \(2\sigma\) of the mean.
        \item Approximately \(99.7\%\) of the data falls within \(3\sigma\) of the mean.
    \end{itemize}




\vs

\textbf{Standardization (Z-Score)} \\
Because every dataset has a different mean and variance, calculating probabilities is difficult. We solve this by transforming any normal distribution into the \textbf{Standard Normal} $Z \sim N(0, 1)$.

    \begin{equation}
    Z=\frac{X-\mu}{\sigma} \sim N(0,1)
    \end{equation}

       \textbf{Why it is important:} In this way, we do not need a different table for every pair \((\mu,\sigma)\). We can always use the table of the standard normal distribution instead.


\begin{figure}[H] 
    \centering\includegraphics[scale=0.42]{111.Standardize Normal Distribution.jpeg}
    \caption{Standardized Normal Distribution}
    \captionsetup{font=footnotesize}
    \label{fig-Standardize Normal Distribution}    
\end{figure}


\textit{Data Science Takeaway:} This exact mathematical operation is called \textit{Standardization.} In Machine Learning, you will constantly use functions like scikit-learn's \texttt{StandardScaler} to apply this formula to your features, ensuring neural networks train efficiently.

\subsubsection{Sampling Distributions (Transformations of Normal)}
In the real world, you never know the true population mean ($\mu$) or variance ($\sigma^2$). You only have a sample. When working with sample statistics, we must use \textit{adapted} versions of the Normal distribution.

\begin{defbox}{1. Chi-Square Distribution \(\chi^2_g\)}
\textbf{What it is:} A Chi-Square random variable is defined as the sum of the squares of \(g\) independent standard normal variables:
\[
X = Z_1^2 + Z_2^2 + \cdots + Z_g^2,
\qquad \text{where } Z_i \sim N(0,1)
\]
Since each term is squared, the variable can only take non-negative values. For this reason, the Chi-Square distribution is defined on \([0,+\infty)\) and is typically right-skewed.

\begin{itemize}
    \item \textbf{Parameter \(g\):} \(g\) is called the number of \textit{degrees of freedom}.

    \item \textbf{Mean:}
    \[
    E[X]=g
    \]

    \item \textbf{Variance:}
    \[
    \mathrm{Var}(X)=2g
    \]

    \item \textbf{Shape:} The distribution is always non-negative. For small \(g\), it is strongly right-skewed; as \(g\) increases, it becomes more symmetric.
    \item \textbf{Important application:} In goodness-of-fit and independence tests, the Chi-Square statistic is often computed as
    \[
    \chi^2=\sum_{i=1}^{k}\frac{(O_i-E_i)^2}{E_i}
    \]
    where \(O_i\) are the observed frequencies and \(E_i\) are the expected frequencies.

\end{itemize}
\end{defbox}




\begin{figure}[H] 
    \centering\includegraphics[scale=0.4]{112.Chi-square and probability table.jpeg}
    \caption{Chi-square distributions for different degrees of freedom and probability table}
    \captionsetup{font=footnotesize}
    \label{fig-Chi-square and probability table}
\end{figure}


The graph shows the density of a Chi-Square random variable. Since the Chi-Square distribution is based on squared quantities, it can only take non-negative values and is typically right-skewed.

\begin{itemize}
    \item \textbf{Interpretation of the graph:} The horizontal axis contains the possible values of the test statistic \(\chi^2\), while the vertical axis represents the density \(f(x)\). The shaded area to the right of a given value \(\chi^2\) represents the probability.   Which is the right-tail probability.
    \[
    P(X \geq \chi^2)
    \]
  
\vs
    \item \textbf{Interpretation of the table:} The Chi-Square table provides critical values for different degrees of freedom and right-tail probabilities. Each row corresponds to a number of degrees of freedom, and each column corresponds to a tail probability level.

    \vs
    \item \textbf{Example:} If the degrees of freedom are \(4\), the value \(9.49\) in the table corresponds to a right-tail probability of \(0.05\). This means:
    \[
    P(X \geq 9.49)=0.05 \qquad \text{for } X \sim \chi^2_4
    \]
    Therefore, if an observed test statistic is greater than \(9.49\), the p-value is smaller than \(0.05\), and this gives evidence against the null hypothesis.

    \vs

    \item \textbf{Typical use:} The Chi-Square distribution is commonly used in goodness-of-fit tests, tests of independence, and inference on population variance.
\end{itemize}

\begin{example}{Practical example of a Chi-Square Goodness-of-Fit test}
Suppose we want to check whether a die is fair.

\begin{itemize}
    \item \textbf{Null hypothesis \(H_0\):} the die is fair, so each face has probability \(1/6\).
    \item \textbf{Alternative hypothesis \(H_1\):} the die is not fair.
\end{itemize}

We roll the die \(60\) times and observe the following frequencies:

\[
O = (5,\,8,\,9,\,12,\,14,\,12)
\]

If the die is fair, the expected frequency for each face is:

\[
E_i=\frac{60}{6}=10
\]

So the expected frequencies are:

\[
E = (10,\,10,\,10,\,10,\,10,\,10)
\]

The Chi-Square test statistic is:

\[
\chi^2 = \sum_{i=1}^{6}\frac{(O_i-E_i)^2}{E_i}
\]

Substituting the values:

\[
\chi^2 =
\frac{(5-10)^2}{10}+
\frac{(8-10)^2}{10}+
\frac{(9-10)^2}{10}+
\frac{(12-10)^2}{10}+
\frac{(14-10)^2}{10}+
\frac{(12-10)^2}{10}
\]

\[
\chi^2 =
\frac{25}{10}+\frac{4}{10}+\frac{1}{10}+\frac{4}{10}+\frac{16}{10}+\frac{4}{10}
= 2.5+0.4+0.1+0.4+1.6+0.4 = 5.4
\]

So the observed test statistic is:

\[
\chi^2_{\text{obs}}=5.4
\]

This value measures how far the observed frequencies are from the expected ones. A small value means that the observed data are close to what we would expect under \(H_0\), while a large value means that the difference is too large to be explained by random variation alone.

\vs


\textbf{Degrees of freedom in the Chi-Square example}
In this die example, there are \(6\) categories, so the degrees of freedom are:

\[
g = 6-1 = 5
\]

\textbf{Why do we subtract 1?}

Because once we know the frequencies of \(5\) categories, the last one is automatically determined by the fact that the total number of observations must be \(60\).

For example, if we already know that the first five observed frequencies are:

\[
5,\;8,\;9,\;12,\;14
\]

their sum is:

\[
5+8+9+12+14=48
\]

Since the total must be \(60\), the last frequency cannot be chosen freely. It must be:

\[
60-48=12
\]

So only \(5\) frequencies are free to vary, and this is why the test has \(5\) degrees of freedom.

\textbf{Interpretation:} Degrees of freedom represent the number of values that can vary independently once the constraint on the total is taken into account.

\end{example}


\vs


\begin{defbox}{2. Student's t-Distribution (\(t_g\))}
The Student's t-distribution is similar to the Normal distribution: it is bell-shaped, symmetric around 0, and defined on all real numbers. However, it has \textbf{heavier tails}, which means that more probability is placed far from the center. 

\vs

This reflects the extra uncertainty that appears when the population variance is unknown and must be estimated from the sample.

\vs
As the sample size increases, the t-distribution becomes closer and closer to the standard Normal distribution.

\begin{itemize}

    \item \textbf{Test statistic:} In the one-sample case, the t-distribution appears in the statistic
    \begin{equation}
    t=\frac{\bar{X}-\mu_0}{S/\sqrt{n}}
    \end{equation}
    where \(\bar{X}\) is the sample mean, \(\mu_0\) is the hypothesized population mean, \(S\) is the sample standard deviation, and \(n\) is the sample size.

    \vs

      \item  \textbf{When to use it:} The t-distribution is used when making inference about a population \textbf{mean}, especially when the sample size is small and the true population variance is unknown.

   \vs
   \item \textbf{Why the tails are heavier:} Since the population standard deviation is not known, we replace it with the sample standard deviation. This introduces additional variability, so the t-distribution gives more weight to extreme values than the Normal distribution.


\end{itemize}
\end{defbox}






\begin{figure}[H] 
    \centering\includegraphics[scale=0.38]{113.t-student and probability table.jpeg}
    \caption{t-Student distributions compared to normal, and probability table}
    \captionsetup{font=footnotesize}
    \label{fig-t-student and probability table}
\end{figure}

   \textbf{Interpretation of the graph:} The curve is centered at \(0\), and the shaded right tail represents a probability of \(0.05\). For \(g=20\) degrees of freedom, the critical value is approximately
    \[
    t_{0.05,20} \approx 1.725
    \]
    meaning that
    \[
    P(T \geq 1.725)=0.05
    \qquad \text{for } T \sim t_{20}
    \]

   \textbf{Interpretation of the table:} Each row corresponds to a number of degrees of freedom, and each column corresponds to a right-tail probability. In the row \(g=20\) and the column \(0.05\), the value \(1.725\) is the critical value that leaves \(5\%\) of the area in the right tail.


   \textbf{Degrees of freedom:} The parameter \(g\) is the number of degrees of freedom. In the one-sample case,
    \[
    g=n-1
    \]
    because one degree of freedom is lost when estimating the sample mean.

\vs

\begin{example}{\textbf{Practical example of a one-sample t-test}}
    
Suppose we want to test whether the average battery life of a device is equal to 10 hours.

\begin{itemize}
    \item \textbf{Null hypothesis \(H_0\):}
    \[
    \mu = 10
    \]

    \item \textbf{Alternative hypothesis \(H_1\):}
    \[
    \mu \neq 10
    \]

    \item \textbf{Sample data:} Assume we collect a sample of \(n=21\) devices, and we compute:
    \[
    \bar{x}=10.8, \qquad s=2.1
    \]

    \item \textbf{Test statistic:}
    \[
    t=\frac{\bar{x}-\mu_0}{s/\sqrt{n}}
    \]

    Substituting the values:
    \[
    t=\frac{10.8-10}{2.1/\sqrt{21}}
    \approx \frac{0.8}{0.458}
    \approx 1.75
    \]

    \item \textbf{Degrees of freedom:}
    \[
    g=n-1=21-1=20
    \]

    \item \textbf{Decision rule:} From the t-table, for \(g=20\) and right-tail probability \(0.05\), the critical value is about
    \[
    1.725
    \]

    Since the observed statistic is
    \[
    t_{\text{obs}} \approx 1.75 > 1.725
    \]
    it falls in the right-tail rejection region at the 5\% level for a one-sided test.

    \item \textbf{Interpretation:} The sample mean is far enough from the hypothesized population mean that it is unlikely to be explained only by random sampling variation.
\end{itemize}
\end{example}

\vspace{0.3cm}

\begin{defbox}{3. F-Distribution \(F_{g,h}\): The Ratio of Two Variances}
\textbf{What it is:} The F-distribution arises as the ratio of two independent Chi-Square variables, each divided by its own degrees of freedom:
\begin{equation}
F=\frac{U/g}{V/h}
\qquad \text{with } U\sim \chi^2_g,\; V\sim \chi^2_h
\end{equation}
It is defined only for positive values and is typically right-skewed.

\begin{itemize}
 

    \item \textbf{Degrees of freedom:} The parameters \(g\) and \(h\) are the degrees of freedom of the numerator and denominator, respectively.

    \item \textbf{When to use it:} The F-distribution is used when comparing two variances, or more generally when comparing explained variability to unexplained variability.

    \item \textbf{Key interpretation:} If \(F\approx 1\), the two sources of variability are similar. If \(F\) is much larger than 1, the numerator variability is substantially larger than the denominator variability.

    \item \textbf{Data Science Use Case:} The F-distribution is the theoretical basis of \textbf{ANOVA} and is also used in regression analysis to test whether a model explains a significant amount of variability.
\end{itemize}
\end{defbox}
\begin{figure}[H] 
    \centering\includegraphics[scale=0.45]{114. F-Snedecor-Fisher and Probability Table.jpeg}
    \caption{F-Snedecor-Fisher distributions and probability table}
    \captionsetup{font=footnotesize}
    \label{fig-F-Snedecor-Fisher and Probability Table}
\end{figure}

 \textbf{Interpretation of the graph:} The horizontal axis contains the possible values of the statistic \(F\), while the vertical axis represents the density. The curve is not symmetric and has a long right tail. The shaded right-tail area represents the probability of observing a value at least as large as the chosen critical value.


 \textbf{Interpretation of the table:} The F-table provides critical values for different pairs of degrees of freedom. One degree of freedom is associated with the numerator and the other with the denominator. In the example shown, the critical value is approximately
    \[
    F \approx 3.18
    \]
    for a specific pair of degrees of freedom and significance level. This means that values larger than \(3.18\) fall in the right-tail rejection region.

    \begin{example}{\textbf{Practical example of the F-distribution: comparing two variances}}

Suppose we want to compare the variability of two production processes.

\begin{itemize}
    \item \textbf{Goal:} Check whether the two populations have the same variance.

    \item \textbf{Null hypothesis \(H_0\):}
    \[
    \sigma_1^2=\sigma_2^2
    \]

    \item \textbf{Alternative hypothesis \(H_1\):}
    \[
    \sigma_1^2 \neq \sigma_2^2
    \]

    \item \textbf{Sample data:} Assume we collect two independent samples:
    \[
    n_1=10, \qquad n_2=12
    \]
    with sample variances
    \[
    s_1^2=18, \qquad s_2^2=6
    \]

    \item \textbf{Test statistic:} We compute the ratio of the two sample variances:
    \[
    F=\frac{s_1^2}{s_2^2}=\frac{18}{6}=3
    \]

    \item \textbf{Degrees of freedom:}
    \[
    g=n_1-1=9, \qquad h=n_2-1=11
    \]

    \item \textbf{Interpretation:} If the two population variances are really equal, then the ratio of the sample variances should be close to 1. A value much larger than 1 suggests that the first population has greater variability than the second one.

    \item \textbf{Decision idea:} We compare the observed value \(F=3\) with the critical value from the F-table for the chosen significance level and the pair of degrees of freedom \((9,11)\). If \(F\) falls in the right-tail rejection region, we reject \(H_0\).
\end{itemize}
\end{example}

